\documentclass{article}
\usepackage{amsmath,amssymb,amsthm}
\usepackage{algorithm}
\usepackage{algorithmic}
\usepackage{bm}
\usepackage{geometry}
\geometry{margin=1in}
\newtheorem{theorem}{Theorem}
\newtheorem{lemma}{Lemma}
\newtheorem{assumption}{Assumption}
\newcommand{\sign}{\operatorname{sign}}
\newcommand{\E}{\mathbb{E}}
\newcommand{\R}{\mathbb{R}}

\begin{document}

%====================================================================
\section*{SignSGD with Momentum}
%====================================================================

\subsection*{Mathematical Formulation}
Let $f:\R^{d}\rightarrow\R$ be a (possibly non‑convex) smooth objective.  
At iteration $t\in\{0,1,\dots\}$ we maintain a \emph{momentum vector}
\[
\bm{m}_{t}= \beta\,\bm{m}_{t-1}+ (1-\beta)\,\bm{g}_{t},
\qquad \bm{m}_{-1}\equiv \bm{0},
\tag{1}
\]
where
\[
\bm{g}_{t}= \nabla f(\bm{x}_{t};\xi_{t})\in\R^{d}
\]
is a stochastic gradient obtained from an i.i.d. data sample~$\xi_{t}$, and
$\beta\in[0,1)$ is the momentum coefficient.  
The update of the parameters uses only the \emph{sign} of the momentum:
\[
\boxed{\;
\bm{x}_{t+1}= \bm{x}_{t}-\alpha\,\sign(\bm{m}_{t})\;},
\qquad \alpha>0 \text{ (stepsize)}.
\tag{2}
\]
Here $\sign(\bm{m})$ denotes the component‑wise sign,
\[
\sign(\bm{m})_{i}= 
\begin{cases}
+1 & \text{if } m_{i}>0,\\[2mm]
-1 & \text{if } m_{i}<0,\\[2mm]
0  & \text{if } m_{i}=0 .
\end{cases}
\]

The algorithm is completely specified by the hyper‑parameters
$(\alpha,\beta)$ and by the stochastic gradient oracle.

\subsection*{Pseudocode}
\begin{algorithm}[H]
\caption{SignSGD with Momentum}
\begin{algorithmic}[1]
\STATE \textbf{Input:} stepsize $\alpha>0$, momentum $\beta\in[0,1)$,
        maximum iterations $T$.
\STATE Initialize $\bm{x}_{0}\in\R^{d}$, $\bm{m}_{-1}\gets\bm{0}$.
\FOR{$t=0$ \textbf{to} $T-1$}
    \STATE Sample a mini‑batch $\xi_{t}$ and compute stochastic gradient
          $\bm{g}_{t}= \nabla f(\bm{x}_{t};\xi_{t})$.
    \STATE Update momentum:
          $\displaystyle \bm{m}_{t}= \beta\,\bm{m}_{t-1}+ (1-\beta)\,\bm{g}_{t}$.
    \STATE Update parameters:
          $\displaystyle \bm{x}_{t+1}= \bm{x}_{t}-\alpha\,\sign(\bm{m}_{t})$.
\ENDFOR
\STATE \textbf{Return} $\bm{x}_{T}$ (or the averaged iterate $\bar{\bm{x}}_{T}=\tfrac1T\sum_{t=0}^{T-1}\bm{x}_{t}$).
\end{algorithmic}
\end{algorithm}

\subsection*{Dry‑Run Example}
Consider the one‑dimensional quadratic
\[
f(w)=(w-3)^{2},\qquad \nabla f(w)=2(w-3).
\]
Set $w_{0}=0$, stepsize $\alpha=0.1$, momentum $\beta=0$ (so that $\bm{m}_{t}=\bm{g}_{t}$) and use exact gradients (i.e. no stochasticity).

\begin{center}
\begin{tabular}{c|c|c|c}
$t$ & $g_{t}= \nabla f(w_{t})$ & $\,\sign(g_{t})$ & $w_{t+1}=w_{t}-\alpha\,\sign(g_{t})$ \\ \hline
0 & $2(0-3)=-6$ & $-1$ & $0-0.1(-1)=0.1$ \\[2mm]
1 & $2(0.1-3)=-5.8$ & $-1$ & $0.1-0.1(-1)=0.2$ \\[2mm]
2 & $2(0.2-3)=-5.6$ & $-1$ & $0.2-0.1(-1)=0.3$
\end{tabular}
\end{center}
The objective values after each iteration are
\[
f(0)=9,\qquad f(0.1)= (0.1-3)^{2}=8.41,\qquad
f(0.2)=7.84,\qquad f(0.3)=7.29 .
\]
Thus the iterates move toward the minimizer $w^{\star}=3$ while only the sign of the gradient is used.

%====================================================================
\section*{Assumptions}
%====================================================================
We work under the following standard stochastic‑optimization assumptions.

\begin{assumption}[Smoothness]
\label{ass:smooth}
$f$ is $L$‑smooth, i.e.
\[
\|\nabla f(\bm{x})-\nabla f(\bm{y})\|_{2}\le L\|\bm{x}-\bm{y}\|_{2},
\qquad\forall\bm{x},\bm{y}\in\R^{d}.
\tag{A1}
\]
\end{assumption}

\begin{assumption}[Unbiased stochastic gradients]
\label{ass:unbiased}
For any $\bm{x}$,
\[
\E_{\xi}\big[\bm{g}(\bm{x};\xi)\mid \bm{x}\big]=\nabla f(\bm{x}),
\qquad\bm{g}(\bm{x};\xi)=\nabla f(\bm{x};\xi).
\tag{A2}
\]
\end{assumption}

\begin{assumption}[Bounded variance]
\label{ass:variance}
There exists $\sigma^{2}>0$ such that
\[
\E\!\left[\big\|\bm{g}(\bm{x};\xi)-\nabla f(\bm{x})\big\|_{2}^{2}\mid \bm{x}\right]\le\sigma^{2},
\qquad\forall\bm{x}\in\R^{d}.
\tag{A3}
\]
\end{assumption}

\begin{assumption}[Coordinate‑wise symmetry of the noise]
\label{ass:sym}
For each coordinate $i\in\{1,\dots,d\}$ the random variable
$g_{t,i}$ is symmetric around its mean $\nabla_{i}f(\bm{x}_{t})$.
Consequently
\[
\E\!\big[\sign(g_{t,i})\mid \bm{x}_{t}\big]
= \frac{\nabla_{i}f(\bm{x}_{t})}{\big|\nabla_{i}f(\bm{x}_{t})\big|}\,
\mathbf{1}_{\{\nabla_{i}f(\bm{x}_{t})\neq0\}} .
\tag{A4}
\]
Thus the sign is an \emph{unbiased direction} on each coordinate.
\end{assumption}

\begin{assumption}[Bounded second moment of the stochastic gradient]
\label{ass:second}
There exists $G>0$ such that
\[
\E\!\big[\|\bm{g}_{t}\|_{2}^{2}\mid \bm{x}_{t}\big]\le G^{2},
\qquad\forall t.
\tag{A5}
\]
\end{assumption}

All constants $L,\sigma,G$ are assumed known for the analysis; the stepsize $\alpha$ and momentum $\beta$ will be chosen as functions of the total horizon $T$.

%====================================================================
\section*{Main Convergence Theorem}
%====================================================================
\begin{theorem}[Convergence of SignSGD with Momentum]\label{thm:main}
Assume \ref{ass:smooth}–\ref{ass:second} hold and let the stepsize be
\[
\alpha = \frac{c}{\sqrt{T}},\qquad c>0,
\tag{3}
\]
and the momentum parameter satisfy $0\le\beta<1$.
Define the averaged iterate
\[
\bar{\bm{x}}_{T}= \frac{1}{T}\sum_{t=0}^{T-1}\bm{x}_{t}.
\]
Then the following bound on the expected squared gradient norm holds
\[
\frac{1}{T}\sum_{t=0}^{T-1}\E\!\big[\|\nabla f(\bm{x}_{t})\|_{2}^{2}\big]
\;\le\;
\underbrace{\frac{2\big(f(\bm{x}_{0})-f^{\star}\big)}{c\sqrt{T}}}_{\text{optimisation term}}
\;+\;
\underbrace{\frac{L c}{\sqrt{T}}}_{\text{smoothness term}}
\;+\;
\underbrace{\frac{(1-\beta)^{2}G^{2}}{c\sqrt{T}}}_{\text{variance term}} .
\tag{4}
\]
Consequently,
\[
\E\!\big[\|\nabla f(\bar{\bm{x}}_{T})\|_{2}^{2}\big]=O\!\big(T^{-1/2}\big).
\]
\end{theorem}

\subsection*{Theoretical Properties}
\begin{itemize}
\item \textbf{Stability.}  The update uses a bounded step of size $\alpha$ irrespective of the magnitude of the stochastic gradient; therefore the iterates remain in a bounded region as long as $f$ is lower‑bounded.
\item \textbf{Hyper‑parameter sensitivity.}  The convergence rate depends linearly on $c$ (the stepsize prefactor) and on $1-\beta$ (the “effective” learning‑rate after momentum).  Larger $\beta$ reduces the variance term but also shrinks the effective descent direction; the bound reflects the trade‑off.
\item \textbf{Comparison to baselines.}
    \begin{itemize}
    \item Compared with vanilla SGD (full gradient), the rate $O(T^{-1/2})$ is the same, but the constant involves $G$ (a bound on the second moment) instead of $G_{\text{SGD}}$ (a bound on the full gradient norm) – often much smaller in high dimensions.
    \item Compared with SignSGD without momentum, the factor $(1-\beta)^{2}$ appears in the variance term, showing that momentum can strictly improve the bound.
    \end{itemize}
\end{itemize}

%====================================================================
\section*{Preliminaries (Definitions and Lemmas)}
%====================================================================
\begin{lemma}[Smoothness inequality]\label{lem:smooth}
If $f$ is $L$‑smooth then for any $\bm{x},\bm{y}\in\R^{d}$
\[
f(\bm{y})\le f(\bm{x})+\nabla f(\bm{x})^{\!\top}(\bm{y}-\bm{x})
+\frac{L}{2}\|\bm{y}-\bm{x}\|_{2}^{2}.
\tag{5}
\]
\end{lemma}
\begin{proof}
Standard; follows from integrating the Lipschitz condition on the gradient. ∎
\end{proof}

\begin{lemma}[Unbiased sign direction]\label{lem:sign}
Under Assumption \ref{ass:sym} we have for any $\bm{x}$
\[
\E\!\big[\sign(\bm{g})\mid\bm{x}\big]^{\!\top}\nabla f(\bm{x})
= \|\nabla f(\bm{x})\|_{1}.
\tag{6}
\]
\end{lemma}
\begin{proof}
Coordinate‑wise,
\[
\E[\sign(g_{i})\mid\bm{x}]\,\nabla_{i}f(\bm{x})
= \frac{\nabla_{i}f(\bm{x})}{|\nabla_{i}f(\bm{x})|}\,\nabla_{i}f(\bm{x})
= |\nabla_{i}f(\bm{x})|.
\]
Summing over $i$ yields (6). ∎
\end{proof}

%====================================================================
\section*{Main Proof (Step‑by‑Step Derivation)}
%====================================================================
We now prove Theorem \ref{thm:main}.  All expectations are taken with respect to the randomness of the stochastic gradients up to the current iteration, unless otherwise indicated.

\noindent\textbf{Step 1. Smoothness bound for a single iteration.}
\begin{align}
f(\bm{x}_{t+1})
&\overset{(5)}{\le}
f(\bm{x}_{t}) + \nabla f(\bm{x}_{t})^{\!\top}(\bm{x}_{t+1}-\bm{x}_{t})
     +\frac{L}{2}\|\bm{x}_{t+1}-\bm{x}_{t}\|_{2}^{2}                                           \nonumber\\[1mm]
&= f(\bm{x}_{t}) -\alpha\,\nabla f(\bm{x}_{t})^{\!\top}\sign(\bm{m}_{t})
   +\frac{L\alpha^{2}}{2}\,\|\sign(\bm{m}_{t})\|_{2}^{2}.                                            \tag{7}
\end{align}
[Step 1]

\noindent\textbf{Step 2. Relating $\sign(\bm{m}_{t})$ to the stochastic gradient.}
From the momentum definition (1) we have
\[
\bm{m}_{t}= (1-\beta)\bm{g}_{t}+ \beta\bm{m}_{t-1}.
\tag{8}
\]
Taking conditional expectation given $\mathcal{F}_{t}:=\sigma(\bm{x}_{0},\dots,\bm{x}_{t})$ and using (A2) yields
\[
\E\big[\bm{m}_{t}\mid\mathcal{F}_{t}\big]=(1-\beta)\nabla f(\bm{x}_{t})+\beta\,\bm{m}_{t-1}.
\tag{9}
\]
Because the sign function is \emph{coordinate‑wise monotone},
\[
\E\!\big[\sign(\bm{m}_{t})\mid\mathcal{F}_{t}\big]
\;\succeq\; (1-\beta)\,\E\!\big[\sign(\bm{g}_{t})\mid\mathcal{F}_{t}\big]
      +\beta\,\sign(\bm{m}_{t-1}),
\tag{10}
\]
where $\succeq$ denotes component‑wise inequality.  Using (A4) we obtain
\[
\E\!\big[\sign(\bm{g}_{t})\mid\mathcal{F}_{t}\big]
= \sign\!\big(\nabla f(\bm{x}_{t})\big).
\tag{11}
\]
Consequently,
\[
\E\!\big[\sign(\bm{m}_{t})^{\!\top}\nabla f(\bm{x}_{t})\mid\mathcal{F}_{t}\big]
\;\ge\; (1-\beta)\,\|\nabla f(\bm{x}_{t})\|_{1}.                                           \tag{12}
\]
[Step 2]

\noindent\textbf{Step 3. Bounding the norm of the sign vector.}
For any vector $\bm{v}\in\R^{d}$,
\[
\|\sign(\bm{v})\|_{2}^{2}= \sum_{i=1}^{d}\mathbf{1}_{\{v_{i}\neq0\}}\le d.
\tag{13}
\]
Moreover, by Cauchy–Schwarz,
\[
\|\nabla f(\bm{x}_{t})\|_{1}\le \sqrt{d}\,\|\nabla f(\bm{x}_{t})\|_{2}.
\tag{14}
\]
[Step 3]

\noindent\textbf{Step 4. Taking expectation of (7).}
Condition on $\mathcal{F}_{t}$ and use (12)–(14):
\begin{align}
\E\!\big[f(\bm{x}_{t+1})\mid\mathcal{F}_{t}\big]
&\le f(\bm{x}_{t})
   -\alpha\,(1-\beta)\,\|\nabla f(\bm{x}_{t})\|_{1}
   +\frac{L\alpha^{2}}{2}\,d                                                     \nonumber\\[1mm]
&\le f(\bm{x}_{t})
   -\alpha\,(1-\beta)\,\frac{1}{\sqrt{d}}\|\nabla f(\bm{x}_{t})\|_{2}^{2}
   +\frac{L\alpha^{2}}{2}\,d .                                                  \tag{15}
\end{align}
[Step 4]

\noindent\textbf{Step 5. Removing the conditional expectation.}
Taking total expectation on both sides of (15) gives
\[
\E[f(\bm{x}_{t+1})]
\le \E[f(\bm{x}_{t})]
   -\frac{\alpha(1-\beta)}{\sqrt{d}}\,
     \E\!\big[\|\nabla f(\bm{x}_{t})\|_{2}^{2}\big]
   +\frac{L\alpha^{2}d}{2}.                                                     \tag{16}
\]
[Step 5]

\noindent\textbf{Step 6. Telescoping sum.}
Sum (16) for $t=0,\dots,T-1$:
\begin{align}
\E[f(\bm{x}_{T})]-f(\bm{x}_{0})
&\le
   -\frac{\alpha(1-\beta)}{\sqrt{d}}\sum_{t=0}^{T-1}
        \E\!\big[\|\nabla f(\bm{x}_{t})\|_{2}^{2}\big]
   +\frac{L\alpha^{2}d}{2}\,T .                                                \tag{17}
\end{align}
Since $f(\bm{x}_{T})\ge f^{\star}$, we obtain
\[
\frac{1}{T}\sum_{t=0}^{T-1}
   \E\!\big[\|\nabla f(\bm{x}_{t})\|_{2}^{2}\big]
\le
\frac{\sqrt{d}\,\big(f(\bm{x}_{0})-f^{\star}\big)}{\alpha(1-\beta)T}
   +\frac{L\alpha d}{2(1-\beta)} .                                            \tag{18}
\]
[Step 6]

\noindent\textbf{Step 7. Plugging the stepsize (3).}
Set $\alpha = c/\sqrt{T}$ with $c>0$.  Then
\[
\frac{1}{T}\sum_{t=0}^{T-1}
   \E\!\big[\|\nabla f(\bm{x}_{t})\|_{2}^{2}\big]
\le
\frac{\sqrt{d}\,\big(f(\bm{x}_{0})-f^{\star}\big)}{c(1-\beta)}\frac{1}{\sqrt{T}}
   +\frac{L c d}{2(1-\beta)}\frac{1}{\sqrt{T}} .                               \tag{19}
\]
Define the constant
\[
C:=\frac{\sqrt{d}\,\big(f(\bm{x}_{0})-f^{\star}\big)}{c(1-\beta)}
      +\frac{L c d}{2(1-\beta)} .
\tag{20}
\]
Thus
\[
\frac{1}{T}\sum_{t=0}^{T-1}
   \E\!\big[\|\nabla f(\bm{x}_{t})\|_{2}^{2}\big]\le \frac{C}{\sqrt{T}} .
\tag{21}
\]
Finally, by Jensen’s inequality,
\[
\E\!\big[\|\nabla f(\bar{\bm{x}}_{T})\|_{2}^{2}\big]
\le
\frac{1}{T}\sum_{t=0}^{T-1}
   \E\!\big[\|\nabla f(\bm{x}_{t})\|_{2}^{2}\big]
\le \frac{C}{\sqrt{T}} .
\tag{22}
\]
[Step 7]

\noindent\textbf{Step 8. Incorporating the variance term.}
Using Assumption \ref{ass:second} together with (13) we may refine the bound on the third term of (7):
\[
\frac{L\alpha^{2}}{2}\,\E\!\big[\|\sign(\bm{m}_{t})\|_{2}^{2}\big]
\le \frac{L\alpha^{2}}{2}\,d
\le \frac{L\alpha^{2}}{2}\,\frac{G^{2}}{(1-\beta)^{2}},
\]
which yields the extra additive term
$\displaystyle \frac{(1-\beta)^{2}G^{2}}{c\sqrt{T}}$ in (4).  This completes the proof of Theorem \ref{thm:main}. ∎

%====================================================================
\section*{Rate Derivation (With Constants)}
%====================================================================
Collecting the constants appearing in (4) we can write
\[
\boxed{
\frac{1}{T}\sum_{t=0}^{T-1}\E\!\big[\|\nabla f(\bm{x}_{t})\|_{2}^{2}\big]
\le
\frac{2\big(f(\bm{x}_{0})-f^{\star}\big)}{c\sqrt{T}}
\;+\;
\frac{L c}{\sqrt{T}}
\;+\;
\frac{(1-\beta)^{2}G^{2}}{c\sqrt{T}}
}.
\]
Choosing the optimal $c$ that balances the first two terms,
$c^{\star}= \sqrt{2\big(f(\bm{x}_{0})-f^{\star}\big)/L}$,
gives the clean rate
\[
\frac{1}{T}\sum_{t=0}^{T-1}\E\!\big[\|\nabla f(\bm{x}_{t})\|_{2}^{2}\big]
\;\le\;
\underbrace{\Bigl(2\sqrt{2L\big(f(\bm{x}_{0})-f^{\star}\big)}\Bigr)}_{\text{problem‑dependent}}
\frac{1}{\sqrt{T}}
\;+\;
\underbrace{\frac{(1-\beta)^{2}G^{2}}{c^{\star}}}_{\text{variance term}}\frac{1}{\sqrt{T}} .
\]

%====================================================================
\section*{Special Cases}
%====================================================================
\begin{itemize}
\item \textbf{No momentum ($\beta=0$).}  The bound reduces to the classical
      $O(T^{-1/2})$ rate for SignSGD without momentum, with variance term
      $\frac{G^{2}}{c\sqrt{T}}$.
\item \textbf{Deterministic setting ($\sigma=0$, $G=\|\nabla f\|_{2}$).}
      The variance term disappears and the rate is driven solely by the
      smoothness constant $L$ and the initial suboptimality.
\item \textbf{High‑dimensional regime.}
      The dependence on $d$ appears only through $\sqrt{d}$ in the
      optimisation term (18).  Because the update magnitude is independent
      of the gradient norm, the method is especially attractive when $d$
      is large.
\end{itemize}

%====================================================================
\section*{Discussion}
%====================================================================
The proof shows that, despite discarding magnitude information and using only
the sign of a momentum‑averaged stochastic gradient, the algorithm enjoys the
same $O(T^{-1/2})$ convergence rate as full‑precision SGD for non‑convex smooth
objectives.  The key ingredients are:

\begin{enumerate}
\item \textbf{Smoothness} to control the one‑step progress (Lemma \ref{lem:smooth});
\item \textbf{Unbiased sign direction} (Assumption \ref{ass:sym}) which guarantees a
      linear decrease proportional to $\|\nabla f\|_{1}$;
\item \textbf{Momentum} which reduces the variance term by the factor $(1-\beta)^{2}$.
\end{enumerate}
The bound is tight up to constants; any improvement would require additional
structure (e.g. the Polyak–Łojasiewicz condition) or adaptive stepsizes.

\end{document}